\documentclass[a4paper,12pt]{article}

\usepackage[utf8]{inputenc}
\usepackage[T1]{fontenc}
\usepackage[turkish]{babel}
\usepackage{geometry}
\usepackage{listings}
\usepackage{xcolor}
\usepackage{booktabs}
\usepackage{array}
\usepackage{hyperref}
\usepackage{tcolorbox}
\usepackage{tikz}
\usepackage{amssymb}

\geometry{margin=2.5cm}

\definecolor{codebg}{RGB}{245,245,245}
\definecolor{codeframe}{RGB}{200,200,200}
\definecolor{keyword}{RGB}{0,0,180}
\definecolor{comment}{RGB}{0,128,0}
\definecolor{string}{RGB}{163,21,21}

\lstdefinestyle{cppstyle}{
    language=C++,
    backgroundcolor=\color{codebg},
    frame=single,
    rulecolor=\color{codeframe},
    basicstyle=\ttfamily\small,
    keywordstyle=\color{keyword}\bfseries,
    commentstyle=\color{comment}\itshape,
    stringstyle=\color{string},
    showstringspaces=false,
    breaklines=true,
    tabsize=4,
    morekeywords={consteval,constexpr,static_assert,decltype,nullptr,auto,requires},
    literate={ç}{{\c{c}}}1 {ş}{{\c{s}}}1 {ğ}{{ğ}}1 {ü}{{\"u}}1 {ö}{{\"o}}1 {ı}{{\i}}1
             {Ç}{{\c{C}}}1 {Ş}{{\c{S}}}1 {Ğ}{{Ğ}}1 {Ü}{{\"U}}1 {Ö}{{\"O}}1 {İ}{{\.I}}1
}
\lstset{style=cppstyle}

\newtcolorbox{notebox}{
    colback=blue!5,
    colframe=blue!40,
    fonttitle=\bfseries,
    title={Önemli Not},
    sharp corners,
    boxrule=0.5pt
}

\title{\textbf{Ders 1 --- Value Categories, Referans Türleri ve \texttt{decltype} Kuralları}}
\author{İleri C++ (C++26 uyumlu)}
\date{24 Mart 2025}

\begin{document}

\maketitle

\noindent\textbf{Kapsam:} Expression value categories, referans bağlama kuralları, \texttt{decltype} tür çıkarımı

\tableofcontents
\newpage

%==============================================================================
\section{İfade (Expression) Temelleri}

Her ifadenin iki temel niteliği vardır:

\begin{table}[h]
\centering
\begin{tabular}{@{}ll@{}}
\toprule
\textbf{Nitelik} & \textbf{Açıklama} \\
\midrule
\textbf{Data Type} (Tür) & İfadenin ürettiği değerin türü (\texttt{int}, \texttt{double}, \texttt{std::string} vb.) \\
\textbf{Value Category} & İfadenin bellekteki konumu ve taşınabilirliğiyle ilgili sınıflandırması \\
\bottomrule
\end{tabular}
\end{table}

\subsection{İfade Türü Hakkında Kritik Kurallar}

\begin{itemize}
    \item İfadelerin türü \textbf{referans türü olamaz}. \texttt{int\& r = x;} bildiriminde \texttt{r} ifadesinin türü \texttt{int}'tir, \texttt{int\&} değil.
    \item İfadelerin türü \textbf{pointer türü olabilir}: \texttt{int* p = \&x;} $\rightarrow$ \texttt{p} ifadesinin türü \texttt{int*}.
    \item \texttt{void} bir türdür; \texttt{void} fonksiyona yapılan çağrı ifadesinin türü \texttt{void}'dir.
    \item \textbf{Integral promotion} kuralı: \texttt{char c; +c} ifadesinin türü \texttt{int}'tir (karakter aritmetiği \texttt{int}'e yükseltilir).
\end{itemize}

\begin{lstlisting}
char c1, c2;
// c1 + c2 ifadesinin turu int'tir (integral promotion)
static_assert(std::is_same_v<decltype(c1 + c2), int>);
\end{lstlisting}

%==============================================================================
\section{Value Categories (Değer Kategorileri)}

\subsection{Birincil (Primary) Değer Kategorileri}

Her ifade \textbf{tam olarak bir} birincil kategoriye aittir:

\begin{table}[h]
\centering
\begin{tabular}{@{}llcc@{}}
\toprule
\textbf{Kategori} & \textbf{Tam Adı} & \textbf{Kimliği Var mı?} & \textbf{Taşınabilir mi?} \\
\midrule
\textbf{lvalue}  & Locator Value  & Evet  & Hayır \\
\textbf{prvalue} & Pure R-Value   & Hayır & Evet  \\
\textbf{xvalue}  & Expiring Value & Evet  & Evet  \\
\bottomrule
\end{tabular}
\end{table}

\subsection{Bileşik (Combined) Değer Kategorileri}

\begin{table}[h]
\centering
\begin{tabular}{@{}lll@{}}
\toprule
\textbf{Kategori} & \textbf{Eşdeğer} & \textbf{Tanım} \\
\midrule
\textbf{glvalue} & lvalue $\cup$ xvalue & Kimliği olan (has identity) tüm ifadeler \\
\textbf{rvalue}  & prvalue $\cup$ xvalue & Taşınabilir (can be moved from) tüm ifadeler \\
\bottomrule
\end{tabular}
\end{table}

\subsection{Venn Şeması}

\begin{center}
\begin{tikzpicture}[
    every node/.style={font=\small},
    box/.style={draw, rounded corners=2pt, minimum width=2.5cm, minimum height=1.5cm, align=center}
]
    % glvalue ellipse
    \draw[thick, blue!60] (-1.5,0) ellipse (3.5cm and 2cm);
    \node[blue!60, font=\bfseries] at (-1.5,2.3) {glvalue};

    % rvalue ellipse
    \draw[thick, red!60] (1.5,0) ellipse (3.5cm and 2cm);
    \node[red!60, font=\bfseries] at (1.5,2.3) {rvalue};

    % lvalue box
    \node[box, fill=blue!10] at (-3,0) {\textbf{lvalue}\\identity\\\textit{!movable}};

    % xvalue box (intersection)
    \node[box, fill=purple!10] at (0,0) {\textbf{xvalue}\\identity\\\textit{movable}};

    % prvalue box
    \node[box, fill=red!10] at (3,0) {\textbf{prvalue}\\!identity\\\textit{movable}};
\end{tikzpicture}
\end{center}

\begin{notebox}
xvalue hem glvalue hem rvalue kümesinin elemanıdır. Bu, xvalue ifadelerin hem kimliğe sahip hem de taşınabilir olduğu anlamına gelir --- move semantiğinin temel taşıdır.
\end{notebox}

\begin{notebox}
Value category bir \textbf{değişkenin} değil, bir \textbf{ifadenin} niteliğidir. \texttt{int\&\& r = 10;} bildiriminde \texttt{r} ifadesinin value category'si \textbf{lvalue}'dur (bir değişken ismi olduğu için), data type'ı ise \texttt{int}'tir.
\end{notebox}

%==============================================================================
\section{Hangi İfade Hangi Kategoride?}

\subsection{lvalue İfadeler}

\begin{table}[h]
\centering
\small
\begin{tabular}{@{}p{4cm}p{4cm}p{6cm}@{}}
\toprule
\textbf{İfade Türü} & \textbf{Örnek} & \textbf{Açıklama} \\
\midrule
Değişken isimleri & \texttt{x}, \texttt{name} & Rvalue ref bile olsa: \texttt{int\&\& r = 5;} $\rightarrow$ \texttt{r} lvalue \\
Fonksiyon isimleri & \texttt{foo} & \texttt{std::move(foo)} bile lvalue üretir \\
String literalleri & \texttt{"hello"} & Dizi gibi davranır $\rightarrow$ lvalue \\
Prefix \texttt{++}/\texttt{-{}-} & \texttt{++a}, \texttt{-{}-a} & Nesnenin kendisini döndürür \\
Atama operatörleri & \texttt{a = b}, \texttt{a += b} & Compound assignment dahil \\
Üye erişimi & \texttt{obj.member} & Nesne lvalue ise erişim de lvalue \\
\texttt{.*} / \texttt{->*} & \texttt{x.*mp} & Member pointer erişimi \\
Dizi erişimi & \texttt{arr[n]} & Dizi lvalue ise \\
Dereference & \texttt{*p} & Her zaman lvalue \\
lvalue ref döndüren fonksiyon & \texttt{foo()} (\texttt{int\& foo();}) & Geri dönüş türü \texttt{T\&} ise \\
lvalue ref'e cast & \texttt{static\_cast<T\&>(expr)} & Hedef tür lvalue referans ise \\
\bottomrule
\end{tabular}
\end{table}

\begin{lstlisting}
#include <string>
#include <utility>

int& get_ref();
void example() {
    int x = 42;
    int&& rr = std::move(x);

    // Hepsi lvalue:
    x;                          // degisken ismi
    rr;                         // rvalue ref degiskeni AMA ifade olarak lvalue!
    "hello world";              // string literal
    ++x;                        // prefix increment
    get_ref();                  // lvalue ref donduren fonksiyon
    *(&x);                      // dereference
}
\end{lstlisting}

\begin{notebox}
\texttt{int\&\& rr = std::move(x);} bildiriminde \texttt{rr}'nin \textbf{data type}'ı \texttt{int}'tir (ifade türü referans olamaz) ve \textbf{value category}'si \textbf{lvalue}'dur (bir isimdir). Tür ile kategoriyi karıştırmak en sık yapılan hatadır.
\end{notebox}

\subsection{prvalue İfadeler}

\begin{table}[h]
\centering
\small
\begin{tabular}{@{}p{5.5cm}p{5.5cm}@{}}
\toprule
\textbf{İfade Türü} & \textbf{Örnek} \\
\midrule
Tüm literaller (string hariç) & \texttt{42}, \texttt{3.14}, \texttt{'A'}, \texttt{true}, \texttt{nullptr} \\
Aritmetik operatörler & \texttt{a + b}, \texttt{a * b} \\
Karşılaştırma operatörleri & \texttt{a == b}, \texttt{a <=> b} \\
Mantıksal operatörler & \texttt{a \&\& b}, \texttt{!a} \\
Postfix \texttt{++}/\texttt{-{}-} & \texttt{a++}, \texttt{a-{}-} \\
Adres operatörü & \texttt{\&x} \\
\texttt{this} pointeri & Üye fonksiyon içinde \texttt{this} \\
Enumerator'ler & \texttt{Color::Red} \\
Lambda ifadeleri & \texttt{[](int x)\{ return x; \}} \\
Geçici nesne oluşturma & \texttt{std::string\{"temp"\}} \\
Referans döndürmeyen fonksiyon çağrısı & \texttt{foo()} (\texttt{int foo();}) \\
Sign operatörü (\texttt{+}, \texttt{-}) & \texttt{+x}, \texttt{-x} \\
\texttt{requires} ifadeleri & \texttt{requires \{ ... \}} (C++20) \\
\bottomrule
\end{tabular}
\end{table}

\begin{lstlisting}
#include <string>

int foo();
void example() {
    int a = 10;

    // Hepsi prvalue:
    42;                         // integer literal
    a + 1;                      // aritmetik
    a > 0;                      // karsilastirma
    &a;                         // adres operatoru
    foo();                      // non-ref donduren fonksiyon
    std::string{"gecici"};      // temporary object
    [](int x){ return x * 2; };// lambda ifadesi
    +a;                         // sign operatoru
}
\end{lstlisting}

\begin{notebox}
Lambda ifadeleri derleyici tarafından closure type türünden \textbf{geçici nesne} (temporary object) olarak ele alınır, bu nedenle her zaman prvalue'dur. \texttt{auto f = []\{\};} bildiriminde \texttt{f} ifadesi lvalue'dur (isimlendirilmiş değişken), ama sağ taraftaki lambda ifadesinin kendisi prvalue'dur.
\end{notebox}

\subsection{xvalue İfadeler}

\begin{table}[h]
\centering
\small
\begin{tabular}{@{}p{5.5cm}p{5.5cm}@{}}
\toprule
\textbf{İfade Türü} & \textbf{Örnek} \\
\midrule
Rvalue ref döndüren fonksiyon çağrısı & \texttt{std::move(x)}, \texttt{std::forward<T>(x)} \\
rvalue nesnenin üye erişimi & \texttt{std::move(obj).member}, \texttt{MyClass\{\}.x} \\
rvalue nesne + \texttt{.*} & \texttt{std::move(obj).*mp} \\
rvalue dizi erişimi & \texttt{std::move(arr)[0]} \\
rvalue ref'e cast & \texttt{static\_cast<T\&\&>(expr)} \\
\bottomrule
\end{tabular}
\end{table}

\begin{lstlisting}
#include <utility>

struct Widget {
    int data;
};

int&& rval_func();

void example() {
    Widget w{42};

    // Hepsi xvalue:
    std::move(w);                  // std::move -> T&& dondurur
    rval_func();                   // rvalue ref donduren fonksiyon
    Widget{}.data;                 // prvalue nesnenin member erisimi
    static_cast<Widget&&>(w);      // rvalue ref'e cast

    // Dizi ornegi:
    using Arr5 = int[5];
    Arr5{1,2,3,4,5}[0];           // rvalue dizinin elemanina erisim -> xvalue
}
\end{lstlisting}

%==============================================================================
\section{Referans Türleri ve Bağlama Kuralları}

\subsection{Üç Referans Kategorisi}

\begin{lstlisting}
std::string& r1  = s;    // (1) lvalue reference
std::string&& r2 = expr; // (2) rvalue reference
// (3) forwarding (universal) reference - yalnizca tur cikarimi baglaminda:
template<typename T>
void func(T&& param);    // T&& burada forwarding reference
\end{lstlisting}

\begin{notebox}
\texttt{T\&\&} ifadesi \textbf{yalnızca tür çıkarımı varsa} forwarding reference'tır. Aksi halde sıradan rvalue reference'tır. Bu ayrım perfect forwarding için kritiktir.
\end{notebox}

\subsection{Bağlama Kuralları Tablosu}

\begin{table}[h]
\centering
\begin{tabular}{@{}lccc@{}}
\toprule
\textbf{Referans Türü} & \textbf{lvalue} & \textbf{prvalue} & \textbf{xvalue} \\
\midrule
\texttt{T\&} (lvalue ref)             & Bağlanır & \textbf{HATA} & \textbf{HATA} \\
\texttt{const T\&} (const lvalue ref) & Bağlanır & Bağlanır      & Bağlanır \\
\texttt{T\&\&} (rvalue ref)           & \textbf{HATA} & Bağlanır & Bağlanır \\
\texttt{const T\&\&} (const rvalue ref) & \textbf{HATA} & Bağlanır & Bağlanır \\
\bottomrule
\end{tabular}
\end{table}

\begin{lstlisting}
#include <string>
#include <utility>

void example() {
    std::string name{"Necati"};

    // lvalue reference - sadece lvalue'ya baglanir:
    std::string& r1 = name;              // OK: lvalue
    // std::string& r2 = std::string{}; // HATA: prvalue
    // std::string& r3 = std::move(name);// HATA: xvalue

    // const lvalue reference - her seye baglanir:
    const std::string& cr1 = name;              // OK: lvalue
    const std::string& cr2 = std::string{};     // OK: prvalue
    const std::string& cr3 = std::move(name);   // OK: xvalue

    // rvalue reference - sadece rvalue'ya baglanir:
    std::string&& rr1 = std::string{};          // OK: prvalue
    std::string&& rr2 = std::move(name);        // OK: xvalue
    // std::string&& rr3 = name;                // HATA: lvalue
}
\end{lstlisting}

\begin{notebox}
\texttt{const T\&\&} pratikte neredeyse hiç kullanılmaz; ancak function overload resolution kurallarını anlamak için varlığını bilmek gerekir.
\end{notebox}

%==============================================================================
\section{Fonksiyon Geri Dönüş Türü ve Value Category İlişkisi}

\begin{table}[h]
\centering
\begin{tabular}{@{}ll@{}}
\toprule
\textbf{Geri Dönüş Türü} & \textbf{Çağrı İfadesinin Value Category'si} \\
\midrule
\texttt{T} (non-reference)     & \textbf{prvalue} \\
\texttt{T\&} (lvalue reference) & \textbf{lvalue}  \\
\texttt{T\&\&} (rvalue reference) & \textbf{xvalue}  \\
\bottomrule
\end{tabular}
\end{table}

\begin{lstlisting}
#include <string>

std::string  by_value();  // cagri -> prvalue
std::string& by_lref();   // cagri -> lvalue
std::string&& by_rref();  // cagri -> xvalue

// std::move'un xvalue uretmesinin sebebi:
// std::move(x) -> geri donus turu T&& -> xvalue
\end{lstlisting}

%==============================================================================
\section{\texttt{decltype} ile Value Category Tespiti}

\texttt{decltype} specifier'ının \textbf{iki farklı çıkarım kuralı} vardır:

\subsection{Kural 1 --- Operand Bir İsim (Identifier) İse}

Çıkarım, ismin \textbf{bildirim türünü} (declared type) birebir verir:

\begin{lstlisting}
int x = 5;
int& r = x;
int&& rr = 10;
const int& cr = x;

struct S { int mx; };
S obj;

static_assert(std::is_same_v<decltype(x),       int>);
static_assert(std::is_same_v<decltype(r),        int&>);
static_assert(std::is_same_v<decltype(rr),       int&&>);
static_assert(std::is_same_v<decltype(cr),       const int&>);
static_assert(std::is_same_v<decltype(obj.mx),   int>);       // uye erisimi de isim
\end{lstlisting}

\subsection{Kural 2 --- Operand Bir İfade (Expression) İse}

Çıkarım, ifadenin \textbf{value category}'sine göre belirlenir:

\begin{table}[h]
\centering
\begin{tabular}{@{}ll@{}}
\toprule
\textbf{Value Category} & \textbf{\texttt{decltype} Ürettiği Tür} \\
\midrule
\textbf{lvalue}  & \texttt{T\&}  \\
\textbf{prvalue} & \texttt{T}    \\
\textbf{xvalue}  & \texttt{T\&\&} \\
\bottomrule
\end{tabular}
\end{table}

\begin{lstlisting}
#include <utility>

int x = 5;

// Parantez eklemek ismi ifadeye donusturur:
static_assert(std::is_same_v<decltype(x),    int>);    // isim -> declared type
static_assert(std::is_same_v<decltype((x)),   int&>);   // ifade + lvalue -> T&

// Value category tespiti:
static_assert(std::is_same_v<decltype(42),              int>);    // prvalue -> T
static_assert(std::is_same_v<decltype(std::move(x)),    int&&>);  // xvalue  -> T&&
static_assert(std::is_same_v<decltype((x)),             int&>);   // lvalue  -> T&
\end{lstlisting}

\begin{notebox}
\texttt{decltype(x)} ile \texttt{decltype((x))} \textbf{farklı sonuç üretir}. Parantez, ismi bir ifadeye dönüştürür. Bu fark, \texttt{decltype(auto)} kullanımında beklenmeyen referans döndürme hatalarına yol açabilir --- dikkat edilmesi gereken kritik bir noktadır.
\end{notebox}

\subsection{\texttt{decltype} ile Value Category Test Fonksiyonu (C++26)}

\begin{lstlisting}
#include <type_traits>
#include <print>

template<typename T>
consteval void print_value_category() {
    if constexpr (std::is_lvalue_reference_v<T>)
        std::println("lvalue");
    else if constexpr (std::is_rvalue_reference_v<T>)
        std::println("xvalue");
    else
        std::println("prvalue");
}

// Kullanim:
// print_value_category<decltype((expr))>();

int main() {
    int x = 42;
    print_value_category<decltype((x))>();              // lvalue
    print_value_category<decltype((std::move(x)))>();   // xvalue
    print_value_category<decltype((42))>();              // prvalue
}
\end{lstlisting}

%==============================================================================
\section{Özet Karar Ağacı}

\begin{center}
\begin{tikzpicture}[
    level distance=2cm,
    sibling distance=4cm,
    every node/.style={font=\small, align=center},
    edge from parent/.style={draw, ->, thick},
    level 1/.style={sibling distance=6cm},
    level 2/.style={sibling distance=3cm}
]
\node {İfadenin value\\category'si nedir?}
    child {node {İsim mi?\\(değişken, fonksiyon,\\string literal)}
        child {node[draw, fill=blue!10, rounded corners] {\textbf{lvalue}\\{\scriptsize (rvalue ref değişkeni\\bile lvalue!)}}}
    }
    child {node {\texttt{std::move()} /\\rvalue ref döndüren\\fonksiyon mu?}
        child {node[draw, fill=purple!10, rounded corners] {\textbf{xvalue}}}
    }
    child {node {Geçici nesne /\\literal / aritmetik /\\lambda mı?}
        child {node[draw, fill=red!10, rounded corners] {\textbf{prvalue}}}
    };
\end{tikzpicture}
\end{center}

\noindent\textbf{Üye erişimi} (\texttt{obj.member}) için:
\begin{itemize}
    \item \texttt{obj} lvalue ise $\rightarrow$ \textbf{lvalue}
    \item \texttt{obj} rvalue ise $\rightarrow$ \textbf{xvalue}
\end{itemize}

%==============================================================================
\section{Low-Latency ve Yüksek Performans İçin Kritik Noktalar}

\begin{notebox}
Move semantiğinin temel motivasyonu gereksiz kopyalamaları ortadan kaldırmaktır. Bir nesne \textbf{xvalue} kategorisindeyse, kaynaklarının ``çalınabileceği'' (move edilebileceği) anlamına gelir. High-frequency trading, oyun motorları ve gerçek zamanlı sistemlerde bu ayrım doğrudan gecikme (latency) farkı yaratır --- özellikle \texttt{std::string}, \texttt{std::vector} gibi heap allocation yapan türlerde.
\end{notebox}

\begin{notebox}
\texttt{const T\&} her value category'ye bağlanabilmesi sayesinde generic kod yazarken güvenli bir ``catch-all'' olarak kullanılır, ancak bu durumda move optimization devre dışı kalır. Performans-kritik kodda overload setinde \texttt{T\&\&} parametresinin bulunması, rvalue argümanlarda kopyalama yerine taşıma yapılmasını sağlar.
\end{notebox}

\end{document}
